\section{Ambiente de Execução e Metodologia}
\label{sec:ambiente}

Todos os experimentos foram conduzidos na servidora \textit{duo} do Departamento de Informática (DINF/UFPR). A Tabela~\ref{tab:cpu} sumariza as especificações do ambiente de hardware, software e os parâmetros operacionais estipulados.

\begin{table}[H]
\centering
\small
\caption{Configuração do ambiente e parâmetros de execução na servidora \textit{duo}.}
\label{tab:cpu}
\begin{tabular}{@{}ll@{}}
\toprule
\textbf{Parâmetro / Componente} & \textbf{Especificação / Valor Adotado} \\ \midrule
Processador & AMD Ryzen 9 5950X (1 soquete, 16 núcleos físicos / 32 threads lógicas) \\
Microarquitetura & Zen 3 (homogênea, sem divisão de núcleos P e E) \\
Frequência / Turbo Boost & Habilitado pelo firmware (estável em 3{,}6--4{,}7 GHz via telemetria contínua) \\
Governador DVFS & Padrão do sistema operacional (sem privilégio root para alteração) \\
Caches & L1d: 32 KiB/núcleo $\cdot$ L2: 512 KiB/núcleo $\cdot$ L3: $2 \times 32$ MiB compartilhada \\
Memória RAM & 125 GiB DDR4 \\
Sistema Operacional & Ubuntu 24.04 LTS (kernel Linux 6.8.0-49-generic) \\
Compilador & \texttt{gcc} 13.3.0 com flags de otimização máxima \texttt{-O3 -fopenmp} \\
OpenMPI & OpenMPI 5.0.x (ambiente Conda), comunicação intranó (\texttt{--mca btl \^{}tcp}) \\
Afinidade de CPU & \texttt{OMP\_PROC\_BIND=close}, \texttt{OMP\_PLACES=cores} (isolamento nos 16 núcleos físicos) \\
Amostragem Estatística & \textbf{20 repetições} por configuração com média e desvio padrão \\ \bottomrule
\end{tabular}
\end{table}

\subsection{Tipagem Dinâmica da Matriz e Recompilação Automática}
\label{sec:tipagem}

A matriz armazena inteiros do tipo \texttt{mtype} (\texttt{macros.h}). Para entradas com dimensão até $65.535$, compila-se com \texttt{uint16\_t} (2 bytes, alvo \texttt{make u16}), reduzindo a pegada de cache pela metade. Para dimensões superiores ($N, M \ge 65.535$), é obrigatório o uso de \texttt{uint32\_t} (\texttt{-DMTYPE\_32}, alvo \texttt{make u32}) para evitar \textit{overflow}. 

O script de automação dos experimentos (\texttt{run\_experiment.py}) gerencia essa transição de forma autônoma: ao detectar que o próximo caso de teste pertence a uma categoria de grande porte ou giga, ele aciona a recompilação imediata dos binários para o modo apropriado antes de iniciar as execuções.
